% =====================================================================
%  Chương 5 — Hạng mục 5: đóng gói triển khai Web App
%  Mã nguồn: app/app.py, app/static/app.js, app/smoke_test.py
% =====================================================================
\chapter{Triển khai Web App}
\label{ch:webapp}

Chương này đóng gói các mô hình thành một Web App chạy cục bộ. Ứng dụng có hai phân hệ ứng với hai tập dữ liệu: dự báo
giá đóng cửa AMZN của phiên kế tiếp, và dự báo xác suất churn của một người dùng KKBox từ chuỗi 31 ngày nghe nhạc. Người
dùng chọn framework (PyTorch hoặc Keras) và một trong bốn mô hình; mặc định là mô hình có val loss nhỏ nhất
(Chương~\ref{ch:keras}).

\section{Kiến trúc}

Máy chủ là Flask, trả JSON qua bốn nhóm endpoint; giao diện là một trang HTML với hai tab, gọi API bằng \texttt{fetch}
và vẽ biểu đồ bằng Chart.js (Hình~\ref{fig:arch}). Ứng dụng không cần dữ liệu gốc của KKBox để chạy. Mọi thứ nó cần
nằm trong thư mục \texttt{models/}: 16 tệp mô hình và \texttt{metadata.json} do notebook 06 ghi, gồm thống kê chuẩn hoá
của hai tập, ngưỡng phân loại chọn trên validation, tên mô hình tốt nhất và chuỗi hành vi của hai người dùng mẫu.

\begin{figure}[H]
  \centering
  \hinh{f5_arch}
  \caption{Kiến trúc Web App}
  \label{fig:arch}
\end{figure}

Bảng~\ref{tab:api} liệt kê các endpoint.

\begin{table}[H]
  \centering\small
  \caption{REST API của Web App}
  \label{tab:api}
  \begin{tabular}{llL{7.2cm}}
    \toprule
    Phương thức & Đường dẫn & Vào $\to$ ra \\
    \midrule
    GET & \texttt{/api/stock/history} & $\to$ 100 phiên gần nhất: ngày, Close, MA20 \\
    POST & \texttt{/api/stock/predict} & \texttt{\{framework, model\}} $\to$ giá dự báo, \% thay đổi, độ trễ \\
    POST & \texttt{/api/churn/predict} & \texttt{\{preset\}} hoặc \texttt{\{sequence\}} ($31\times8$), \texttt{framework},
                                        \texttt{model} $\to$ xác suất, mức rủi ro, ngưỡng, độ trễ \\
    GET & \texttt{/api/meta} & $\to$ danh sách mô hình, mô hình tốt nhất, người dùng mẫu \\
    GET & \texttt{/} & $\to$ trang giao diện \\
    \bottomrule
  \end{tabular}
\end{table}

\section{Nạp mô hình}

Mỗi mô hình chỉ được nạp một lần, ở lần gọi đầu tiên, rồi giữ trong bộ nhớ nhờ \texttt{lru\_cache}. Hàm nạp trả về một
hàm dự báo cùng giao diện cho cả hai framework, nên phần xử lý yêu cầu không cần biết mô hình đến từ đâu. Mô hình Keras
chạy qua \texttt{tf.function} như lúc đo độ trễ ở Chương~\ref{ch:keras}, và mỗi mô hình được chạy thử một lần ngay khi
nạp để độ trễ trả về cho người dùng là độ trễ lúc ổn định. TensorFlow chỉ được import khi có yêu cầu dùng mô hình Keras.

\maNguon{app__load_model}{Nạp mô hình một lần, cùng giao diện cho PyTorch và Keras}

\section{Hai phân hệ dự báo}

Phân hệ giá lấy 30 phiên mới nhất trong \texttt{AMZN.csv}, áp đúng phép biến đổi của Chương~\ref{ch:dulieu} (giá tương
đối, log Volume, z-score bằng thống kê của train), đưa qua mô hình rồi quy đổi return dự báo ra giá USD. Hàm
\texttt{latest\_window} dùng chung hai hàm biến đổi với \texttt{make\_windows} của lúc huấn luyện, nên dữ liệu phục vụ
không thể lệch khỏi dữ liệu huấn luyện.

\maNguon{app__stock_predict}{Endpoint dự báo giá phiên kế tiếp}

Phân hệ churn nhận hoặc tên một người dùng mẫu, hoặc một ma trận $31 \times 8$ bất kỳ. Đầu vào sai shape bị từ chối với mã
400 và thông báo lỗi rõ ràng, không đưa vào mô hình. Xác suất được so với ngưỡng của đúng mô hình đó để xếp mức rủi ro:
cao khi vượt ngưỡng, trung bình khi vượt nửa ngưỡng, còn lại là thấp.

\maNguon{app__churn_predict}{Endpoint dự báo churn}

\section{Giao diện}

Mỗi tab có một cột điều khiển và một biểu đồ. Khi đổi framework, danh sách mô hình được dựng lại và mô hình tốt nhất
trên validation của framework đó được chọn sẵn.

\maNguonJS{initChurn}{Tab churn: gọi API và vẽ chuỗi hành vi bằng Chart.js}

Hình~\ref{fig:ui-stock} và~\ref{fig:ui-churn} vẽ lại hai tab bằng TikZ thay cho ảnh chụp màn hình. Mọi con số và đường
cong trong hai hình là phản hồi thật của API, do \texttt{smoke\_test.py} ghi lại: biểu đồ giá đọc 100 phiên mà
\texttt{/api/stock/history} trả về, còn biểu đồ churn đọc đúng chuỗi mà \texttt{/api/churn/predict} trả về cho người dùng
mẫu.

\begin{figure}[H]
  \centering
  \hinh{f5_ui_stock}
  \caption{Tab dự báo giá AMZN với mô hình PyTorch tốt nhất trên validation}
  \label{fig:ui-stock}
\end{figure}

Với dữ liệu kết thúc ngày \SL{app/stock/torch/last_date}, mô hình \SL{app/stock/torch/model_label} của PyTorch dự báo
phiên kế tiếp đóng cửa ở \SL{app/stock/torch/pred_close} USD, mô hình \SL{app/stock/keras/model_label} của Keras dự báo
\SL{app/stock/keras/pred_close} USD, so với giá hiện tại \SL{app/stock/torch/last_close} USD. Cả hai chỉ lệch vài phần
nghìn so với giá hiện tại, đúng với điều Chương~\ref{ch:keras} đã phân tích: mô hình dự báo return gần 0.

\begin{figure}[H]
  \centering
  \hinh{f5_ui_churn}
  \caption{Tab dự báo churn với một người dùng thật đã churn trong tập test}
  \label{fig:ui-churn}
\end{figure}

Hai người dùng mẫu là người thật trong tập test: một người đã churn mà mô hình tin chắc nhất, và một người ở lại mà mô hình
cho xác suất churn thấp nhất. Người churn trong Hình~\ref{fig:ui-churn} nghe gần như mỗi ngày trong 9 ngày đầu tháng rồi
ngừng hẳn từ ngày 10, dạng cực đoan của quỹ đạo nhóm churn ở Chương~\ref{ch:dulieu}. LSTM của PyTorch cho xác suất
\SL{app/churn/churn_torch/probPct}\,\% với người này và \SL{app/churn/loyal_torch/probPct}\,\% với người ở lại; GRU của
Keras cho lần lượt \SL{app/churn/churn_keras/probPct}\,\% và \SL{app/churn/loyal_keras/probPct}\,\%.

\section{Kiểm chứng}

\texttt{smoke\_test.py} dùng test client của Flask, không cần mở cổng mạng. Nó gọi lịch sử giá, dự báo giá với cả hai
framework và dự báo churn cho hai người dùng mẫu với cả hai framework, tổng cộng \SL{app/n_calls} lượt. Mỗi phản hồi phải
có mã 200 và đủ các khoá cần thiết; một chuỗi sai shape phải bị từ chối với mã \SL{app/bad_shape_status}. Kịch bản dừng
ngay khi có một điều kiện không thoả, và lượt chạy cuối qua hết.

\maNguon{smoke_test__main}{Kiểm chứng ba endpoint và ghi phản hồi thật cho báo cáo}

Chạy ứng dụng bằng \texttt{python app/app.py} rồi mở \texttt{http://localhost:8080}. Độ trễ API trả về là thời gian của riêng lời gọi mô hình:
\SL{app/churn/churn_torch/latency_ms} ms với LSTM của PyTorch và \SL{app/churn/churn_keras/latency_ms} ms với GRU của
Keras, cùng cỡ với số đo ở Bảng~\ref{tab:bench-kkbox}. Trong ứng dụng, PyTorch chạy 1 luồng còn TensorFlow dùng số luồng
mặc định, và mỗi con số chỉ là một lần gọi, nên chúng để tham khảo chứ không thay cho phép đo có kiểm soát của
Chương~\ref{ch:keras}.
